\documentclass[11pt]{article}

\usepackage[margin=1in]{geometry}
\usepackage{amsmath,amssymb,amsthm,amsfonts}
\usepackage{mathtools}
\usepackage{bm}
\usepackage{hyperref}
\usepackage{enumitem}
\usepackage{microtype}
\usepackage[T1]{fontenc}
\usepackage[utf8]{inputenc}

\title{Neuromorphic AGI, Multiplicity Theory, and the Universal Multiplicity Equation:\\
From Descriptive PDE to Operator Generator and Executable Simulator}
\author{Citizen Gardens \\ \\ The Foundation of Multiplicity}
\date{April 3, 2026}

\begin{document}
\maketitle

\begin{abstract}
This report synthesizes developments around Multiplicity Theory, Neuromorphic Artificial General Intelligence (AGI), and the Universal Multiplicity Equation (UME). The original neuromorphic AGI preprint presents a time-dependent multiplicity equation as a unifying dynamical law for cognitive state evolution, integrating prime-based encoding, recursive feedback, tensor interactions, diffusion, and stochasticity. Here we upgrade that formulation in three directions: (1) we re-cast the UME from a descriptive partial differential equation into a generator of a semigroup on a multiplicity Hilbert space; (2) we formalize prime indices as labels for irreducible operator families acting on this space; and (3) we provide a concrete PyTorch-style architecture (the MultiplicityCell) together with an experiment design to test a prime-factor transfer law across composite tasks. The report concludes with recommendations for next-step experiments and validation pathways.
\end{abstract}

\tableofcontents

\section{Executive Summary}

\subsection{Starting point: neuromorphic UME}

The neuromorphic AGI preprint introduces a Universal Multiplicity Equation of the form
\begin{equation}
\frac{\partial \psi_k(t)}{\partial t}
= \alpha_k(t)\psi_k
+ \beta_k(t)\int_0^t I_k(\tau)\,d\tau
+ \gamma_k(t)\sum_{j,l} T_{kjl}\psi_j\psi_l
+ \lambda_k(t)\nabla^2\psi_k
+ \eta_k(t)\psi_k^n
+ \xi_k(t),
\label{eq:UME-coordinate}
\end{equation}
with explicit dimensional analysis and characteristic scales ensuring that each term matches the dimensions of the left-hand side.\footnote{All references to the neuromorphic AGI preprint and its equations are drawn from the attached document.} This equation is presented as a framework for cognitive state evolution in neuromorphic AGI, combining:
\begin{itemize}[nosep]
\item Growth or decay via coefficients $\alpha_k$,
\item Memory integrals via $\beta_k$ and $I_k$,
\item Tensor-mediated coupling via $T_{kjl}$,
\item Diffusion via $\lambda_k\nabla^2$,
\item Nonlinear self-interaction via $\eta_k\psi_k^n$,
\item Stochasticity via $\xi_k$.
\end{itemize}

The preprint further incorporates quantum-inspired structure via prime-encoded superpositions, tensorial entanglement, and geometric feedback terms using Berry curvature and the Fubini--Study metric.

\subsection{Conceptual upgrade}

We reinterpret \eqref{eq:UME-coordinate} not merely as a PDE written in components, but as the coordinate representation of an operator equation
\begin{equation}
\frac{d\Psi}{dt} = \mathcal{L}_{\mathrm{mult}}[\Psi],
\end{equation}
where $\Psi$ is a state in a (possibly infinite-dimensional) multiplicity Hilbert space $\mathcal{H}$, and $\mathcal{L}_{\mathrm{mult}}$ is a generator of a strongly continuous semigroup (or group, when unitary) incorporating distinct operator contributions.

A schematic operator decomposition is
\begin{equation}
\mathcal{L}_{\mathrm{mult}}
=
\Xi(t) + \Lambda_m(t) + \mathcal{T} + \mathcal{D} + \mathcal{N} + \mathcal{S},
\end{equation}
where:
\begin{itemize}[nosep]
\item $\Xi(t)$ encodes prime evolution,
\item $\Lambda_m(t)$ is a multiplicity-scalar operator,
\item $\mathcal{T}$ is a tensor coupling operator,
\item $\mathcal{D}$ is a diffusion operator,
\item $\mathcal{N}$ is a nonlinear operator,
\item $\mathcal{S}$ is a stochastic term.
\end{itemize}

Prime indices are elevated from static labels to indices of irreducible operator families $\{\hat{O}_p\}$. Composite cognition is modeled as operator words $W = \prod_p \hat{O}_p^{k_p}$ acting on $\Psi$. Ethics is upgraded from a linear update rule to a geometric constraint manifold $\mathcal{E}\subset\mathcal{H}$ defined by a fairness functional $F[\Psi]\ge \varepsilon$.

\subsection{Executable upgrade}

We then construct a concrete MultiplicityCell in PyTorch that discretizes the dynamics
\[
\Psi_{t+1} = \Psi_t + \Delta t\,\mathcal{L}_{\mathrm{mult}}(\Psi_t)
\]
over a discrete multiplicity axis representing prime-encoded basis states. The cell includes:
\begin{enumerate}[nosep]
\item A diagonal prime-evolution operator,
\item A nonlinear multiplicity-scalar operator,
\item A tensor-like mixing operator across prime indices,
\item A discrete Laplacian as diffusion operator,
\item Nonlinear self-interaction,
\item Optional ethical projection and geometric regularization via Fubini--Study curvature penalties.
\end{enumerate}

On top of this cell, we design a multi-task experiment to test a prime-factor transfer law across tasks, where transfer performance is predicted to scale with the normalized overlap of their activated prime sets.

\section{From Coordinate UME to Operator Generator}

\subsection{Multiplicity Hilbert space}

We introduce a multiplicity Hilbert space $\mathcal{H}$ which, in one simple realization, could be
\[
\mathcal{H} \cong \ell^2(\mathbb{P}) \otimes \mathcal{K},
\]
where:
\begin{itemize}[nosep]
\item $\ell^2(\mathbb{P})$ is the square-summable space over primes, with canonical basis $\{\lvert p_k\rangle\}$,
\item $\mathcal{K}$ is a feature space (for example $L^2(\Omega)$ over a spatial domain $\Omega$, or $\mathbb{C}^d$ for finite feature vectors).
\end{itemize}

A general state can then be written as
\[
\Psi = \sum_{k} \psi_k \otimes \lvert p_k\rangle,
\]
with $\psi_k \in \mathcal{K}$ corresponding to the component associated with prime $p_k$.

\subsection{Operator decomposition of $\mathcal{L}_{\mathrm{mult}}$}

We define $\mathcal{L}_{\mathrm{mult}}$ as an operator (possibly time-dependent) on $\mathcal{H}$, decomposed schematically as
\begin{equation}
\mathcal{L}_{\mathrm{mult}}
=
\Xi(t) + \Lambda_m(t) + \mathcal{T} + \mathcal{D} + \mathcal{N} + \mathcal{S}.
\end{equation}

\paragraph{Prime evolution operator $\Xi(t)$}

We let
\[
\Xi(t) = \sum_{k} \alpha_k(t)\, \mathbb{I} \otimes \lvert p_k\rangle\langle p_k\rvert,
\]
which acts diagonally on the prime index, matching the growth or decay term $\alpha_k(t)\psi_k$.

\paragraph{Multiplicity-scalar operator $\Lambda_m(t)$}

This is a nonlinear or linear operator acting within each prime block, for example
\[
\Lambda_m(t)[\Psi] = \sum_k \Lambda^{(k)}_m(t)\psi_k \otimes \lvert p_k\rangle,
\]
where $\Lambda_m^{(k)}(t)$ may be realized as a small neural network or MLP on the feature space $\mathcal{K}$.

\paragraph{Tensor coupling operator $\mathcal{T}$}

We interpret the $T_{kjl}$ tensor as defining an operator-valued quadratic form:
\begin{equation}
\mathcal{T}[\Psi] = \sum_{k,j,l} T_{kjl}\,(\psi_j,\psi_l)\otimes \lvert p_k\rangle,
\end{equation}
where $(\psi_j,\psi_l)$ is a suitable bilinear combination (for example pointwise product or another bilinear map). In a low-rank factorization,
\[
T_{kjl} = \sum_{r=1}^R U_{kr}V_{jr}W_{lr},
\]
leading to efficient tensor network implementations.

\paragraph{Diffusion operator $\mathcal{D}$}

We let
\[
\mathcal{D} = \sum_k \lambda_k(t) \Delta \otimes \lvert p_k\rangle\langle p_k\rvert,
\]
where $\Delta$ is a Laplacian on $\mathcal{K}$ (spatial, graph, or otherwise). On a discrete prime axis, we may also consider a Laplacian in the prime index itself.

\paragraph{Nonlinear operator $\mathcal{N}$}

The nonlinear self-interaction term $\eta_k(t)\psi_k^n$ is captured as
\[
\mathcal{N}[\Psi] = \sum_k \eta_k(t)\, f_n(\psi_k)\otimes \lvert p_k\rangle,
\]
where $f_n$ denotes a power or general nonlinear map.

\paragraph{Stochastic term $\mathcal{S}$}

The stochastic term $\xi_k(t)$ is modeled as noise:
\[
\mathcal{S}[\Psi,t] = \Xi_{\mathrm{noise}}(t),
\]
with $\Xi_{\mathrm{noise}}(t)$ a stochastic process in $\mathcal{H}$.

\subsection{Prime operators and composite cognition}

Instead of treating primes as mere labels, we introduce a family of operators $\{\hat{O}_p\}$ on $\mathcal{H}$, each associated with a prime $p$. Composite cognition is then an operator word
\[
W = \prod_p \hat{O}_p^{k_p},
\]
acting on $\Psi$. A task or cognitive transformation can be represented as such an operator word, or as a word in a noncommutative algebra generated by the $\hat{O}_p$.

\subsection{Ethics as constraint manifold}

We upgrade the discrete ethical update rule
\[
M(t+1) = \alpha M(t) + \beta R(t)
\]
to a geometric constraint:
\begin{equation}
\Psi(t)\in\mathcal{E}\subset\mathcal{H},
\end{equation}
with
\[
\mathcal{E} = \{\Psi\in\mathcal{H} \mid F[\Psi] \geq \varepsilon\},
\]
where $F$ is a fairness or lawfulness functional. The dynamics are then restricted to $\mathcal{E}$ either by
\begin{itemize}[nosep]
\item projected evolution, $\Psi_{t+1} = \Pi_{\mathcal{E}}\big(\Psi_t + \Delta t\,\mathcal{L}_{\mathrm{mult}}(\Psi_t)\big)$, or
\item penalization in a variational functional, for example,
\[
\mathcal{F}[\Psi] = E[\Psi] + R[\Psi] + C[\Psi] + \lambda_{\mathcal{E}}\max(0,\varepsilon - F[\Psi]).
\]
\end{itemize}

\subsection{Quantum geometry via Fubini--Study distance}

On the projective Hilbert space $\mathbb{P}(\mathcal{H})$, the Fubini--Study distance between normalized states $\Psi_1,\Psi_2$ is
\[
d_{\mathrm{FS}}(\Psi_1,\Psi_2)
=
\arccos\big(|\langle \Psi_1,\Psi_2\rangle|\big).
\]
A curvature penalty can be introduced over trajectories $\Psi(t)$ by integrating or summing discrete Fubini--Study distances between successive states.

\section{Variational Learning and Prime-Factor Transfer Law}

\subsection{Variational viewpoint}

We treat learning as a constrained variational problem:
\begin{equation}
\Psi^\ast = \arg\min_{\Psi \in \mathcal{E}} \mathcal{F}[\Psi],
\end{equation}
where
\begin{equation}
\mathcal{F} = E[\Psi] + R[\Psi] + C[\Psi],
\end{equation}
and $E[\Psi]$ is a stability or energy functional, $R[\Psi]$ is task loss, and $C[\Psi]$ is a curvature or geometry penalty.

\subsection{Prime-factor transfer law}

Let each task $T_i$ be associated with a subset $P_i$ of prime indices. Define
\[
\mathrm{Overlap}(i,j) = \frac{|P_i\cap P_j|}{|P_i\cup P_j|}.
\]
The prime-factor transfer law posits
\begin{equation}
\mathrm{Transfer}(T_i\to T_j)
\approx
a\,\mathrm{Overlap}(i,j) + b
\end{equation}
for constants $a>0$, $b$, where transfer is measured via zero-shot performance, few-shot sample efficiency, or reduction in loss after a fixed adaptation budget.

\section{Discrete Simulator: MultiplicityCell}

\subsection{Architecture}

We discretize the UME dynamics over time with step size $\Delta t$, on a finite-dimensional approximation of $\mathcal{H}$ given by tensors of shape
\[
\text{batch} \times \text{feature-dim} \times \text{multiplicity-dim}.
\]
Here, the multiplicity dimension indexes prime basis states.

A PyTorch-style definition of one time-step of the dynamics is:

\begin{verbatim}
import torch
import torch.nn as nn
import torch.nn.functional as F

class MultiplicityCell(nn.Module):
    def __init__(self, feature_dim, multiplicity_dim,
                 hidden_dim, delta_t=0.1, nonlinear_order=3):
        super().__init__()
        self.feature_dim = feature_dim
        self.multiplicity_dim = multiplicity_dim
        self.hidden_dim = hidden_dim
        self.delta_t = delta_t
        self.nonlinear_order = nonlinear_order

        self.alpha = nn.Parameter(torch.randn(multiplicity_dim) * 0.1)

        self.scalar_net = nn.Sequential(
            nn.Linear(feature_dim, hidden_dim),
            nn.ReLU(),
            nn.Linear(hidden_dim, feature_dim)
        )

        self.mixing_matrix = nn.Parameter(
            torch.randn(multiplicity_dim, multiplicity_dim) * 0.1
        )
        self.coupling_proj = nn.Linear(feature_dim, feature_dim)

        self.lambda_diff = nn.Parameter(torch.randn(multiplicity_dim) * 0.1)
        self.register_buffer("laplacian",
                             self._compute_laplacian(multiplicity_dim))

        self.nonlinear_eta = nn.Parameter(torch.randn(multiplicity_dim) * 0.01)

        self.ethical_projection = None

    def _compute_laplacian(self, n):
        L = torch.diag(torch.full((n,), -2.0))
        for i in range(n - 1):
            L[i, i + 1] = 1.0
            L[i + 1, i] = 1.0
        return L

    def forward(self, psi, external_input=None):
        batch_size = psi.size(0)

        term1 = self.alpha.view(1, 1, -1) * psi

        psi_flat = psi.permute(0, 2, 1).reshape(-1, self.feature_dim)
        scalar_out = self.scalar_net(psi_flat)
        term2 = scalar_out.reshape(batch_size, self.multiplicity_dim,
                                   self.feature_dim).permute(0, 2, 1)

        mixed = torch.einsum('kj, bfj -> bfk', self.mixing_matrix, psi)
        term3 = self.coupling_proj(mixed.permute(0, 2, 1)).permute(0, 2, 1)

        lap_psi = torch.einsum('kj, bfj -> bfk', self.laplacian, psi)
        term4 = self.lambda_diff.view(1, 1, -1) * lap_psi

        term5 = self.nonlinear_eta.view(1, 1, -1) * (psi ** self.nonlinear_order)

        if external_input is None:
            noise = torch.randn_like(psi) * 0.01
        else:
            noise = external_input

        L_mult = term1 + term2 + term3 + term4 + term5 + noise

        psi_next = psi + self.delta_t * L_mult

        if self.ethical_projection is not None:
            psi_next = self.ethical_projection(psi_next)

        return psi_next
\end{verbatim}

\subsection{Geometric regularization}

A discrete Fubini--Study proxy and curvature penalty can be implemented as:

\begin{verbatim}
def compute_fs_distance(psi1, psi2):
    b = psi1.size(0)
    v1 = psi1.view(b, -1)
    v2 = psi2.view(b, -1)
    v1 = F.normalize(v1, dim=1)
    v2 = F.normalize(v2, dim=1)
    overlap = torch.abs(torch.sum(v1 * v2, dim=1))
    overlap = torch.clamp(overlap, -0.999, 0.999)
    return torch.acos(overlap)

def curvature_penalty(psi_seq):
    total = 0.0
    for t in range(len(psi_seq) - 1):
        dist = compute_fs_distance(psi_seq[t], psi_seq[t + 1])
        total += dist.mean()
    return total
\end{verbatim}

\section{Outlook}

We have:
\begin{itemize}[nosep]
\item Recast the neuromorphic UME as an operator-generator framework on multiplicity space.
\item Clarified primes as indices of operator families and compositional modes.
\item Upgraded the ethical component to a constraint manifold.
\item Instantiated a concrete MultiplicityCell architecture that discretizes the UME.
\item Sketched an experiment to test a prime-factor transfer law across tasks with controlled prime-support overlap.
\end{itemize}

Implementing and running the associated experiments is the next step to turn this conceptual unification into empirically grounded, falsifiable claims.

\nocite{*}
\bibliographystyle{plainnat}
\bibliography{references} 
\end{document}